\part{Introducción}

Este documento está orientado a guiarte en el proceso de generar una hoja de vida profesional en {\lmr\LaTeX}.
Esta plantilla ha sido diseñada bajo el formato \texttt{.dtx} ({\lmr Documented \LaTeX}), lo que significa que el código fuente de la plantilla y este manual de usuario coexisten en un único archivo modular e independiente.
Si compilas este documento se obtiene esta guía de uso y estilo.

El objetivo principal de utilizar esta clase es separar por completo el contenido de la forma.
Mientras te concentras en redactar tus logros y experiencia de manera clara y precisa, {\lmr\LaTeX} se encarga de gestionar la consistencia visual, la jerarquía tipográfica y la optimización del espacio.
Esto garantiza un documento elegante para un evaluador humano y perfectamente legible para los sistemas automatizados de selección de personal (ATS).



\section{Estructura del proyecto y flujo de trabajo}

Esta plantilla se compone estrictamente de dos archivos fundamentales que trabajan en conjunto:

\begin{itemize}
    \item \textbf{Archivo de instalación (\texttt{.ins}):} Es el encargado de generar el motor de diseño. Al compilarlo, se extrae automáticamente el archivo de clase (\texttt{.cls}), el cual contiene toda la lógica, macros y herramientas de formato de la plantilla.
    \item \textbf{Archivo documentado (\texttt{.dtx}):} Es el archivo que estás leyendo en este momento. Su función es doble: por un lado, contiene el código fuente de la plantilla y, por el otro, genera este manual de usuario para explicarte cómo utilizar cada comando.
\end{itemize}



\section{Instrucciones para Empezar}

Para comenzar a construir tu hoja de vida, debes seguir este orden cronológico en tu entorno de desarrollo o terminal:

\begin{enumerate}
    \item \textbf{Generar la documentación:} Para generar el manual de documentación de la clase debes ejecutar el comando:
    \begin{verbatim}
lualatex cv.dtx
    \end{verbatim}

    \item \textbf{Generar la clase:} Compila el archivo \texttt{.ins} utilizando:
    \begin{verbatim}
lualatex cv.ins
    \end{verbatim}
    Esto crea el archivo \texttt{cv.cls} en tu directorio de trabajo.

    \item \textbf{Descargar las fuentes:} Antes de poder compilar el documento, es necesario descargar las fuentes que se utilizan:
    \begin{itemize}
        \item En Windows, ejecute el archivo \texttt{fonts.ps1}. Muy probablemente Windows te pida permisos para ejecutarlo.
        \item En Linux o MacOS, ejecute el archivo \texttt{fonts.sh}, para poder ejecutarlo necesitará \texttt{curl} y \texttt{jq}.
    \end{itemize}

    \item \textbf{Escribir la Hoja de Vida:} Crea un archivo nuevo con extensión \texttt{.tex} (por ejemplo, \texttt{cv.tex}) en la misma carpeta donde se encuentra el archivo \texttt{cv.cls} generado. En la primera línea del documento, invoca la clase utilizando:
    \begin{verbatim}
\documentclass{cv}
    \end{verbatim}
\end{enumerate}

A partir de este punto, puedes utilizar las siguientes partes de este manual para conocer los comandos específicos que esta plantilla te ofrece para dar formato a tus datos personales, educación, experiencia laboral y habilidades.



\section{Cómo leer esta guía}

Este manual está organizado de forma secuencial, siguiendo el mismo orden lógico en el que construyes una hoja de vida.
Comenzando por la configuración global, el encabezado con tus datos personales, y avanzando hacia las secciones.

A lo largo de las secciones, encontrarás bloques técnicos donde se definen los nuevos comandos utilizando la nomenclatura estándar de la documentación de {\lmr\LaTeX}:
\begin{itemize}
    \item Los comandos nuevos aparecen resaltados en el margen izquierdo para su rápida localización.
    \item La sintaxis te indicará los argumentos requeridos entre llaves \marg{argumento} y los opcionales (si los hay) entre corchetes \oarg{argumento}.
\end{itemize}

Te recomiendo mantener este manual abierto como referencia mientras redactas tu documento principal \texttt{.tex}, avanzando sección por sección.
